
\medskip
\phantomsection\label{cor:sga3-VIB-2.8}%
\noindent\textbf{Corollary 2.8.---}%
\footnote{\label{note:sga3-VIB-21}Editor's note: we have changed
the presentation by separating the assertions concerning the
``quasi-finite'' case from those concerning the ``unramified''
case.}
\emph{Let \(u:G\to H\) be a morphism of \(S\)-groups. Suppose
that, for every \(s\in S\), \(G_s\) is locally of finite type over
\(\kappa(s)\).}
\begin{enumerate}
\item[\textup{a)}] \emph{If \(u\) is locally of finite type, the
  following conditions are equivalent:}
  \begin{enumerate}
  \item[\textup{(i)}] \emph{\(u\) is locally quasi-finite;}
  \item[\textup{\SGAInlineRefTarget{cor:sga3-VIB-2.8:item:ii}(ii)}] \emph{for every \(s\in S\), the morphism
    \(u_s:G_s\to H_s\) is locally quasi-finite;}
  \item[\textup{\SGAInlineRefTarget{cor:sga3-VIB-2.8:item:iii}(iii)}] \emph{\(\operatorname{Ker}u\) is locally
    quasi-finite over \(S\);}
  \item[\textup{(iv)}] \emph{the fibers of
    \(\operatorname{Ker}u\) are discrete.}
  \end{enumerate}
\item[\textup{b)}] \emph{If \(u\) is locally of finite
  presentation, the following conditions are equivalent:}
  \begin{enumerate}
  \item[\textup{(v)}] \emph{\(u\) is unramified;}
  \item[\textup{(vi)}] \emph{for every \(s\in S\), the morphism
    \(u_s:G_s\to H_s\) is unramified;}
  \item[\textup{\SGAInlineRefTarget{cor:sga3-VIB-2.8:item:vii}(vii)}] \emph{\(\operatorname{Ker}u\) is
    unramified;}
  \item[\textup{(viii)}] \emph{the identity section
    \(S\to\operatorname{Ker}u\) is an open immersion.}
  \end{enumerate}
\end{enumerate}

The equivalences
\[
  \textup{(i)}\Longleftrightarrow\textup{(ii)}
  \qquad\text{and}\qquad
  \textup{(v)}\Longleftrightarrow\textup{(vi)}
\]
follow from \hyperref[concl:sga3-VIB-2.5-ii]{2.5}\textup{\SGARefLink{cor:sga3-VIB-2.8:item:ii}{(ii)}}, and Reminder~\hyperref[rem:sga3-VIB-2.8.1]{2.8.1} below shows that
\[
  \textup{(i)}\Longrightarrow\textup{(iii)}
  \qquad\text{and}\qquad
  \textup{(v)}\Longrightarrow\textup{(vii)}.
\]

For every \(s\in S\), denote by \(e_s\) the identity element of
\(H_s\). Then \textup{\SGARefLink{cor:sga3-VIB-2.8:item:iii}{(iii)}} (respectively, \textup{\SGARefLink{cor:sga3-VIB-2.8:item:vii}{(vii)}})
implies that, for every \(s\in S\),
\[
  (\operatorname{Ker}u)_s
    =\operatorname{Ker}u_s
    =u_s^{-1}(e_s)
\]
is locally quasi-finite (respectively, unramified) over
\(\kappa(s)=\kappa(e_s)\). It follows that \(u_s\) is quasi-finite
(respectively, unramified) at the identity point of \(G_s\), and
hence, by \SGARefLink{prop:sga3-VIB-1.3}{1.3}, that \(u_s\) is locally quasi-finite (respectively,
unramified). Thus
\[
  \textup{(iii)}\Longrightarrow\textup{(ii)}
  \qquad\text{and}\qquad
  \textup{(vii)}\Longrightarrow\textup{(vi)}.
\]
Finally,
\[
  \textup{(ii)}\Longleftrightarrow\textup{(iv)}
  \quad\text{by 1.4.1,}\qquad
  \textup{(vii)}\Longleftrightarrow\textup{(viii)}
  \quad\text{by 2.7.}
\]

\medskip
\phantomsection\label{rem:sga3-VIB-2.8.1}%
\noindent\textbf{Reminder \SGARefLink{rem:sga3-VIB-2.8.1}{2.8.1}.---}
Recall (cf.\ \SGARefLink{sga3:i:definition:2-3-6-1}{Expos\'e~I, 2.3.6.1}) that, given a morphism
\(u:G\to H\) of \(S\)-groups, the \emph{kernel of \(u\)}, denoted
\(\operatorname{Ker}u\), is the group subfunctor of \(G\) defined,
for every morphism \(f:T\to S\), by
\[
  (\operatorname{Ker}u)(T)
    =\{a\in G(T)\mid u\circ a=\epsilon_H\circ f\}.
\]
By loc.\ cit.\ (or EGA~I, 4.4.1), this functor is represented by
the \(S\)-group
\[
  G\times_H S=u^{-1}\bigl(\epsilon_H(S)\bigr),
\]
which is denoted simply by \(\operatorname{Ker}u\). In particular,
the structural morphism
\[
  \operatorname{Ker}u\longrightarrow S
\]
is obtained from \(u\) by base change.

\medskip
\phantomsection\label{lem:sga3-VIB-2.9}%
\noindent\textbf{Lemma 2.9.---}
\emph{Let \(\pi:X\to S\) be a morphism admitting an \(S\)-section
\(\epsilon\).}
\begin{enumerate}
\item[\textup{(i)}] \emph{If \(\pi\) is injective, then it is
  integral.}%
  \begingroup
  \renewcommand{\thefootnote}{\ensuremath{(\ast)}}%
  \footnote[1]{\label{note:sga3-VIB-2.9-star}This is also a
  special case of EGA~IV\(_4\), 18.12.11, since \(\pi\) is
  plainly a universal homeomorphism.}%
  \endgroup
\item[\textup{(ii)}] \emph{If \(\pi\) is locally of finite type
  and, for every \(s\in S\), \(\pi_s\) is an isomorphism, then
  \(\pi\) is an isomorphism.}%
  \begingroup
  \renewcommand{\thefootnote}{\ensuremath{(\ast\ast)}}%
  \footnote[2]{\label{note:sga3-VIB-2.9-starstar}This is also an
  immediate consequence of EGA~IV\(_4\), 18.12.6.}%
  \endgroup
\end{enumerate}

First observe that, by \hyperref[lem:sga3-VIB-2.9.1]{Lemma~2.9.1} below, \(\pi\), being a
homeomorphism, is an affine morphism.

If \(\pi\) is injective, then \(\epsilon\) is surjective. Since
\(\epsilon\) is a surjective immersion, \(\epsilon(S)\) is
isomorphic to the closed subscheme of \(X\) defined by a nil ideal
\(\mathcal I\) of \(\mathcal O_X\). Since \(\epsilon\) is an
\(S\)-section of \(\pi\), there is a direct-sum decomposition
\[
  \mathcal O_X=\mathcal O_S\oplus\mathcal I
\]
of \(\mathcal O_S\)-modules. Since \(\mathcal I\) is a nil ideal
of \(\mathcal O_X\), it is plainly integral over \(\mathcal O_S\).
Thus \(\mathcal O_X\) is integral over \(\mathcal O_S\), and
\(\pi\) is integral.

Now suppose that \(\pi\) is locally of finite type. Since
\(\pi\circ\epsilon=\operatorname{id}_S\), the section \(\epsilon\)
is locally of finite presentation
(cf.\ EGA~IV\(_1\), 1.4.3(v)); hence \(\mathcal I\) is an ideal of
finite type in \(\mathcal O_X\) (EGA~IV\(_1\), 1.4.1). For every
\(s\in S\), one has
\[
  \mathcal O_{X_s}
    =\kappa(s)\oplus
      \mathcal I\otimes_{\mathcal O_S}\kappa(s).
\]
By hypothesis, \(\epsilon_s\) is an isomorphism, so
\[
  \mathcal I\otimes_{\mathcal O_S}\kappa(s)=0
\]
for every \(s\in S\), and therefore, a fortiori,
\[
  \mathcal I\otimes_{\mathcal O_X}\kappa(x)=0
\]
for every \(x\in X\). Nakayama's lemma now gives
\(\mathcal I=0\), and consequently \(\pi\) is an isomorphism.

\medskip
\phantomsection\label{lem:sga3-VIB-2.9.1}%
\noindent\textbf{Lemma 2.9.1.---}
\emph{Let \(f:X\to Y\) be a morphism of schemes that is a
homeomorphism. Then \(f\) is an affine morphism.}%
\begingroup
\renewcommand{\thefootnote}{\ensuremath{(\ast\ast\ast)}}%
\footnote[3]{\label{note:sga3-VIB-2.9.1-starstarstar}Cf.\
EGA~IV\(_4\), 18.12.7.1 for a slightly more general result, proved
by the same argument.}%
\endgroup

It is enough to show that every point \(y\in Y\) has an open
neighborhood \(W\) such that the restriction of \(f\) over \(W\)
is affine. Let \(y\in Y\), and let \(V\) be an affine open
neighborhood of \(y\) in \(Y\). Set
\[
  V'=f^{-1}(V).
\]
Then \(V'\) is an open neighborhood of \(x=f^{-1}(y)\) in \(X\).
There is an affine open neighborhood \(W'\) of \(x\) in \(X\)
contained in \(V'\). Set
\[
  W=f(W').
\]
Then \(W\) is an open neighborhood of \(y\) in \(Y\), contained
in the affine scheme \(V\), and therefore \(W\) is separated.
Since \(W'\) is affine, the restriction of \(f\) over \(W\) is an
affine morphism (EGA~II, 1.2.3).

\medskip
\phantomsection\label{cor:sga3-VIB-2.10}%
\noindent\textbf{Corollary 2.10.---}
\emph{Let \(G\) be an \(S\)-group locally of finite type. Suppose
that, for every \(s\in S\), \(G_s\) is the identity
\(\kappa(s)\)-group. Then \(G\) is the identity \(S\)-group.}

More generally:

\medskip
\phantomsection\label{cor:sga3-VIB-2.11}%
\noindent\textbf{Corollary 2.11.---}
\emph{Let \(f:X\to Y\) be an \(S\)-morphism locally of finite
type. In order that \(f\) be a monomorphism, it is necessary and
sufficient that \(f_s\) be a monomorphism for every \(s\in S\).}

The condition is plainly necessary; we prove that it is
sufficient. If \(f_s\) is a monomorphism for every \(s\in S\),
then, a fortiori, \(f_y\) is a monomorphism for every \(y\in Y\).
We may therefore assume that \(Y=S\).

By EGA~I, 5.3.8, to prove that \(f\) is a monomorphism it is enough
to show that
\[
  \Delta_f:X\longrightarrow X\times_SX
\]
is an isomorphism, or, equivalently, that the first projection
\[
  p:X\times_SX\longrightarrow X
\]
is an isomorphism. If \(f_s\) is a monomorphism, EGA~I, 5.3.8
likewise shows that the first projection
\[
  X_s\times_{\kappa(s)}X_s\longrightarrow X_s,
\]
which is identified with \(p_s\), is an isomorphism. But \(p\)
has the \(S\)-section \(\Delta_f\); hence \hyperref[lem:sga3-VIB-2.9]{Lemma~2.9} shows that, if
\(p_s\) is an isomorphism for every \(s\in S\), then \(p\) itself
is an isomorphism.

\medskip
\phantomsection\label{cor:sga3-VIB-2.12}%
\noindent\textbf{Corollary 2.12.---}
\emph{Let \(G\) be an \(S\)-scheme admitting an \(S\)-section
\(\epsilon\), and suppose that the structural morphism
\(\pi:G\to S\) is closed. Let \(s\in S\) be such that \(\pi\) is
of finite presentation at \(\epsilon(s)\), and such that}
\[
  \epsilon_s:\operatorname{Spec}\bigl(\kappa(s)\bigr)
    \longrightarrow G_s
\]
\emph{is an isomorphism (or, equivalently, such that}
\[
  \pi_s:G_s\longrightarrow
    \operatorname{Spec}\kappa(s)
\]
\emph{is an isomorphism). Then there exists an open neighborhood
\(U\) of \(s\) in \(S\) such that}
\[
  \epsilon|_U:U\longrightarrow G\times_SU
\]
\emph{is an isomorphism.}

Let \(V\) be the set of points of \(G\) at which \(\pi\) is
unramified. This set is known to be open, and it contains \(\epsilon(s)\). Hence
\[
  W=\epsilon^{-1}(V)
\]
is an open subset of \(S\) containing \(s\), and \(\pi\) is
unramified at \(\epsilon(t)\) for every \(t\in W\). Since \(\pi\)
is closed, the same is true of \(\pi|_W\), and we may therefore
suppose that \(W=S\).

By \hyperref[lem:sga3-VIB-2.7.1]{2.7.1},
\[
  X=G-\epsilon(S)
\]
is then a closed subset of \(G\) that does not meet \(G_s\).
Since \(\pi\) is closed,
\[
  F=\pi(X)
\]
is a closed subset of \(S\) that does not contain \(s\). Set
\[
  U=S-F.
\]
Then \(U\) is an open subset of \(S\), and \(\epsilon|_U\) is an
isomorphism from \(U\) onto \(G|_U\).
